% =====================================================================
%  Chapter 13 — Step 7
% =====================================================================

\chapter{Step Seven: Averaging over the Hecke Characters}
\label{ch:averaging}

With the period independent of the character, the argument can average. But
averaging over Hecke characters is not free: the set of characters is a
discrete group, the function being averaged is not constant, and the average
of a non-constant function over a group is not the function's value --- it is
smaller, by a variance term. That variance is exactly what has to be
controlled, and the remark's items 11 and 19 are where the method gets its
tool.

The remark's item 11:

\begin{quote}
\small
``Iwasawa tower via analogy to pigeonhole on $\mathbb{F}_{q^n}/\mathbb{F}_q$.''
\end{quote}

and item 19:

\begin{quote}
\small
``precise choice of the Iwasawa tower, specifically using Behrend!!''
\end{quote}

Behrend's name is the clue. Behrend's 1940 construction is a pigeonhole
argument that finds, inside a large multiplicative group, a subgroup on which
many elements have small order --- and the point of the construction is a
\emph{variance bound}: the spread of a function over a translate of a Behrend
set is small.

\section{Behrend's construction}Behrend's construction, in its classical form, finds a set $T \subseteq
\{1,\dots,N\}$ with no three-term arithmetic progression of small
coefficients, and of size at least
\begin{equation}
N \exp\!\left(-C\sqrt{\log N \log(K+1)}\right),
\label{eq:behrendbound}
\end{equation}
for an absolute constant $C$.

The development states this exactly, and the statement is worth quoting
because it shows the construction being formalised as a theorem about finite
cyclic groups rather than as an analytic heuristic.

\begin{seamnote}{\code{Formalism.BehrendSets.Statement} (excerpt)}
{\footnotesize
\begin{leanquote}
def BehrendFree (K : Nat) (T : Finset Nat) : Prop :=
forall alpha beta gamma : Int, alpha != 0 -> beta != 0 ->
    gamma != 0 -> alpha + beta + gamma = 0 ->
|alpha| <= (K : Int) -> |beta| <= (K : Int) -> |gamma| <= (K
      : Int) ->
    forall x1 in T, forall x2 in T, forall x3 in T,
alpha * (x1 : Int) + beta * (x2 : Int) + gamma * (x3 : Int)
        = 0 -> x1 = x2 and x2 = x3

def BehrendBound : Prop := exists C : Real, 0 < C and
  forall K N : Nat, 1 <= K -> 1 <= N ->
exists T : Finset Nat, T subseteq range N and BehrenFree K T
      and
(N : Real) * Real.exp (-C * Real.sqrt (Real.log N * Real.log
        ((K : Real) + 1))) <= (T.card : Real)

def ThickSet : Prop := exists C : Real, 0 < C and
  forall n L : Nat, 1 <= L -> 8 * L <= n ->
exists T : Finset Nat, (forall t in T, 8 * L * t < n) and
      BehrenFree (2 * L) T and
(n : Real) / (8 * L) * Real.exp (-C * Real.sqrt (Real.log n
        * Real.log (2 * (L : Real) + 1)))
        <= (T.card : Real)
\end{leanquote}
}
\end{seamnote}

Read \code{BehrendFree}: it says that no three distinct elements of $T$ satisfy
a non-trivial linear relation $\alpha x_1 + \beta x_2 + \gamma x_3 = 0$ with
$|\alpha|, |\beta|, |\gamma| \le K$. That is the generalised
three-term-progression-free condition, and it is the right generalisation
because what is needed is not progressions but short relations.

And \code{ThickSet} is the multi-scale version: for each $L$ with $8L \le n$,
a Behrend-free set of the right shape and the right cardinality. The constant
$8$ is the shape of the Behrend argument itself.

\section{Why a Behrend set buys a variance bound}The analytic point is this. Suppose one wants to average a function over
translates of a set $T$, and the translates are indexed by a group. If $T$ has
no short additive relations among its elements, then the translates
``look independent'', and the average of a function over the translates is
close to the average over the whole group, with an error controlled by the
number of short relations --- which is zero.

That is exactly what the development's variance statements record. In
\code{Formalism/BehrendSets/Statement.lean}:

\begin{seamnote}{The variance statements}
{\footnotesize
\begin{leanquote}
def SumSq : Prop := forall M Y : Nat, 2 <= M ->
sum l in (range Y).filter (PmOne M), Real.log l / (l :
    Real)^2
    <= 40 * Real.log M / (M : Real)^2

def SumLin : Prop := forall M Y : Nat, 2 <= M ->
sum l in (range (Y + 1)).filter (PmOne M), Real.log l / (l :
    Real)
    <= 8 * (1 + Real.log Y)^2 / (M : Real)

def ShortOrdLarge : Prop := forall (s Lw l : Nat) (hl :
  Nat.Coprime l (2^(s+2))), 1 <= Lw ->
8 * Lw <= 2^s -> 2 <= l -> orderOf (frob s l hl) <= 4 * Lw
    -> 2^s <= Lw * (l + 1)

def ShortOrdSq : Prop := forall (s Lw Y : Nat), 1 <= Lw -> 8
  * Lw <= 2^s ->
  sum l in (range Y).filter
(fun l => 2 <= l and exists hl : Nat.Coprime l (2^(s+2)),
        orderOf (frob s l hl) <= 4 * Lw),
Real.log l / (l : Real)^2 <= 80 * (Lw : Real)^2 * Real.log
      ((2 : Real)^s) / ((2 : Real)^s)^2
\end{leanquote}
}
\end{seamnote}

These are the quantitative statements. \code{SumSq} and \code{SumLin} bound the
sum of $\log l / l^2$ and $\log l / l$ over a set of $l$ with a given
property, in terms of the modulus $M$ alone --- with absolute constants $40$
and $8$. \code{ShortOrdLarge} and \code{ShortOrdSq} do the same for the
elements $l$ of small order in the subgroup.

The arithmetic here is the classical one. $\sum_l \log l / l^2 \le
40\log M/M^2$ over a set of $l$ with $l \equiv \pm 1 \pmod M$ is what one
needs to bound the logarithm of a product; and \code{OrdIff} records the
order criterion
\begin{equation}
\operatorname{ord}(l) \le 2^j \iff l \equiv \pm 1 \pmod{2^{s+2-j}},
\end{equation}
which is the structure of $(\mathbb{Z}/2^{s+2})^\times$.

\section{Where the group comes from: the Iwasawa tower}The remark's analogy is ``pigeonhole on $\mathbb{F}_{q^n}/\mathbb{F}_q$'', and
the development's group is not $\mathbb{F}_q^\times$ but
$(\mathbb{Z}/2^{s+2})^\times$:

\begin{seamnote}{\code{Formalism.BehrendSets.Statement}}
{\footnotesize
\begin{leanquote}
abbrev OrdGrp (s : Nat) : Type :=
(ZMod (2^(s+2)))^x / Subgroup.zpowers (-1 : (ZMod
    (2^(s+2)))^x)

def frob (s l : Nat) (hl : Nat.Coprime l (2^(s+2))) : OrdGrp
  s :=
  QuotientGroup.mk (ZMod.unitOfCoprime l hl)

def OrdCard : Prop := forall s : Nat, Nat.card (OrdGrp s) =
  2^s and IsCyclic (OrdGrp s)
\end{leanquote}
}
\end{seamnote}

So: take the unit group of $\mathbb{Z}/2^{s+2}$, quotient by the two-element
subgroup generated by $-1$, and get a cyclic group of order $2^s$. The element
\code{frob} is the Frobenius --- the image of $l$ under the $l$-th power map ---
and \code{OrdIff} is the order criterion.

\begin{quote}
\textbf{The connection to the tower.} The point of the construction is now
visible: one wants a group in which the Frobenius elements have small order,
because the order of the Frobenius is what controls the growth of the
corresponding $L$-function. Behrend's construction finds, by pigeonholing, a
large set of elements whose orders are all bounded by $4L_w$, and the
variance bounds \code{ShortOrdSq} and \code{SumSq} then control what those
elements contribute.

In the remark's own phrase, this is the Iwasawa tower ``by analogy to
pigeonhole on $\mathbb{F}_{q^n}/\mathbb{F}_q$'': one wants a nested sequence
of groups whose sizes grow and whose Frobenius elements stay small, which is
precisely what an Iwasawa tower provides and what Behrend's pigeonhole
provides in the finite setting.
\end{quote}

This is where the remark's footnote comes in: ``This is the sort of trick that
intuitively to me is obviously going to be unnecessary.'' The author thinks
the Behrend construction is not actually needed, because in the number-field
setting the tower can be chosen directly. The development contains it anyway.

\section{The equidistribution: the real payload}The final statement of the Behrend subtree is the one that does the work, and
it is the direct formalisation of the remark's item 15 --- ``Glue an integral
structure across Hecke characters for a `gcd', extra averaging''.

\begin{seamnote}{\code{TranslateMean}}
{\footnotesize
\begin{leanquote}
def TranslateMean : Prop := forall (n : Nat) [NeZero n] (T :
  Finset Nat), (forall t in T, t < n) ->
forall {iota : Type} (I : Finset iota) (pos : iota -> ZMod
    n),
sum u : ZMod n, (I.filter (fun i => pos i in thick n T
      u)).card = I.card * T.card
\end{leanquote}
}
\end{seamnote}

Read it. Fix $n$, fix a Behrend set $T \subseteq \{1,\dots,n-1\}$. For each
$u \in \mathbb{Z}/n$, let $\mathrm{thick}(n, T, u)$ be the translate of $T$
by $u$. Now take any finite set $I$ of positions, each with a residue
$\mathrm{pos}(i)$ modulo $n$. Then
\begin{equation}
\sum_{u \in \mathbb{Z}/n} \bigl|\{i \in I : \mathrm{pos}(i) \in \mathrm{thick}(n,T,u)\}\bigr|
\;=\; |I| \cdot |T|.
\label{eq:translatemean}
\end{equation}

\begin{quote}
\textbf{What this says.} Every element of $I$ lies in exactly $|T|$ of the
$n$ translates --- on average, and the identity is exact. The translates
``cover each position equally many times''. This is the equidistribution
that makes the averaging legitimate: whatever function one is averaging, the
average over the translates of a Behrend set is the average over the whole
group, up to an error controlled by the short relations --- and there are none.
\end{quote}

This is the formal statement of the "gcd" and the "extra averaging" of item 15.
It is also why the Behrend set is the right object rather than an arbitrary
one: an arbitrary set of size $|T|$ does not satisfy \eqref{eq:translatemean};
a Behrend set does, and the size of the Behrend set is what makes the error
term small.

The companion statements \code{WindowTwo} and \code{OrbitOne} say that within
a window of length $2L$ at most two elements of the translate occur
(\code{WindowTwo}), and that the orbit map is injective (\code{OrbitOne}) ---
which together are what make the translates genuinely non-overlapping in the
controlled sense.

\section{How this plugs back into the chain}\begin{center}
\begin{tabular}{@{}l>{\RaggedRight\arraybackslash\hyphenpenalty=10000\exhyphenpenalty=10000}p{0.68\textwidth}@{}}
\toprule
\textbf{Remark item} & \textbf{What it provides} \\
\midrule
11, 19 & The group with small-order Frobenius elements, and the quantitative
bounds on what they contribute. \\
15 & The equidistribution \eqref{eq:translatemean}, which makes the average
over the Hecke characters legitimate. \\
8 & The characters $\mu_j$, now sum-able because their period is
character-independent (Chapter~\ref{ch:seesaw}). \\
14 & The budget arithmetic, which now has to pay for the averaging. \\
\bottomrule
\end{tabular}
\end{center}

The last row is the one to keep in mind. Averaging over the Hecke characters
costs something, and the cost lands in the budget of Chapter~\ref{ch:budget}.
That is precisely why the remark's item 14 is dated the same day as the
averaging items, and why the two-thirds appears where it does.

\begin{keytakeaways}
\item Averaging over Hecke characters incurs a variance, and the tool for
controlling it is Behrend's pigeonhole construction: a large set with no
short linear relations among its elements.
\item The development formalises this as \code{BehrendFree}, \code{BehrendBound}
and \code{ThickSet}, together with the quantitative variance bounds
\code{SumSq}, \code{SumLin}, \code{ShortOrdLarge} and \code{ShortOrdSq}, with
absolute constants.
\item The group is $(\mathbb{Z}/2^{s+2})^\times / \langle -1\rangle$, cyclic of
order $2^s$, and the relevant element is the Frobenius, whose order is what
controls the growth.
\item The payload is \code{TranslateMean}: each position lies in exactly
$|T|$ of the $n$ translates. That equidistribution is what makes the average
over the Hecke characters legitimate, and it is the formal content of the
remark's ``gcd'' and ``extra averaging''.
\end{keytakeaways}
